\scanpage{14}
égale à $1$, ce qui équivaut donc au

\uline{Corollaire} Dans $\Gamma_{\mathbb Q}=\Gamma_{\overline{\mathbb Q}_0/\mathbb Q}$, le centralisateur de tout sous-groupe ouvert est réduit à $(1)$.

OPS le sous-groupe ouvert \inserted{$\Gamma'$} invariant, Il est bien connu (?) que son centre est réduit à $1$
\margin{à vérifier ?}
donc si $\unsure{Z|\mathcal Z}$ est son centralisateur dans $\Gamma_{\mathbb Q}$, l'hom $\unsure{Z|\mathcal Z}\longrightarrow\Gamma_{\mathbb Q}/\Gamma'$ est injectif donc $\unsure{Z|\mathcal Z}$ est fini. Mais on sait que les seuls \unsure{éléments|élémts} $\neq1$ de $\Gamma_{\mathbb Q}$ d'ordre fini sont les conjugués de $\tau$, conjugaison complexe. Mais le centralisateur de $\tau$ dans $\Gamma_{\overline{\mathbb Q}}$ est réduit à $(1,\tau)$, donc ne peut contenir $\Gamma'$, donc $\tau\notin\unsure{Z|\mathcal Z}$, donc $\unsure{Z|\mathcal Z}=(1)$.

Revenant à $\unsure{\mathcal E|\epsilon}\subset E_{\bar K/K}$, on voit donc que son centralisateur $\unsure{Z|\mathcal Z}$ dans $E_{\bar K/K}$ a une image dans $\Gamma$ réduite à $\{1\}$, donc $\unsure{Z|\mathcal Z}\subset\pi_{\bar K/K}$. Soit $\pi'\subset\pi=\pi_{\bar K/K}$ la trace de $\unsure{\mathcal E'|\epsilon'}$ sur $\pi$, c'est un sous-groupe ouvert de $\pi$, et on est ramené à
\margin{13}
